\subsection{Archimedes 公理}

本节所讨论的实直线的拓扑性质都是下述定理的推论：

定理 1. 若 \(x\) 与 \(y\) 是任意两个 \(>0\) 的实数，则存在整数 \(n > 0\) 使得 \(y < nx\) 。

存在两个有理数 \(p / q\) 与 \(r / s\) 使得 \(0 < p / q < x\) 且 \(y < r / s\) ，因为开区间 \(]0, x[\) 与 \(]y, \rightarrow[\) 非空（§ 1 第 4 目，命题 4）；取 \(n\) 使得 \(nps > qr\) ，就有 \(y < nx\) 。

注. 实数理论的公理化构造将在第五章 § 2 中给出，在那里定理 1 作为一个公理出现；关于这个公理的更多细节，可参看第四章的历史注记。

